%===============================%
%     IDENTITAS KARYA TULIS     %
%===============================%

% --- Judul ---
%
% Gunakan \judul di bawah ini untuk mengganti judul seperti:
%  - Lembar Jawaban Tutorial Online
%  - Tugas Tutorial Online 1
% Menjadi judul apa pun yang Anda inginkan
%
%\newcommand{\judul}{Judul Tulisan Anda}
%\newcommand{\judulSebaris}{\begingroup\def\\{ }\judul\endgroup}

% --- Keterangan Maksud atau Tujuan Penulisan Karya Tulis ---
%
% Gunakan \keteranganTujuanKarya di bawah ini untuk menulis 
% pernyataan tujuan dari tulisan yang Anda susun ini.
%
% [Contoh] "Diajukan sebagai salah satu syarat untuk memperoleh
%          gelar Sarjana Sains pada Program Studi \programStudi."
%
%\newcommand{\tujuanKarya}{Diajukan sebagai salah satu syarat penuntasan Tugas \tugasKe\ \mataKuliah.}

% --- Tanggal dan Tahun ---
%
% Gunakan \today untuk tanggal lengkap saat ini
% Gunakan \the\year untuk tahun ini
%
\newcommand{\tanggalLengkap}{\today}
\newcommand{\tahun}{\the\year}

% --- Tempat dan Tanggal Lengkap ---
\newcommand{\tempatTanggal}{\daerahMahasiswa, \tanggalLengkap}



%=======================================%
%     IDENTITAS SUBJEK PEMBELAJARAN     %
%=======================================%

% --- Skema Pembelajaran/Tutorial ---
%
% Perintah dua bahasa yang tersedia:
%
%  - \TutorialOnline
%    Tutorial Online <--> Online Tutorial
%
%  - \TutorialWebinar
%    Tutorial Webinar <--> Webinar Tutorial
%
%  - \TutorialTatapMuka
%    Tutorial Tatap Muka <--> Face-to-Face Tutorial
%
\newcommand{\skemaPembelajaran}{\TutorialOnline}

% --- Sesi atau Tugas ---
%
% Bila menggunakan \tugasKe, info \sesiKe akan diabaikan untuk
% halaman sampul, header, dan sebagainya
%
\newcommand{\sesiKe}{1}
%\newcommand{\tugasKe}{1}

% --- Mata Kuliah dan Kode Kelas ---
\newcommand{\kodeMataKuliah}{STSI9999}
\newcommand{\namaMataKuliah}{Pengendalian Berbasis Desktop}

\newcommand{\kodeKelas}{256}



%=============================%
%     IDENTITAS MAHASISWA     %
%=============================%===============================%
% Keterangan Fungsi Variabel
% [Kosong]  : Semua
% <-- ENG   : Dokumen Berbahasa Inggris
% <-- LJM   : Lembar Jawaban
% <-- BT    : Buku Tugas
% <-- KTI   : Karya Tulis Ilmiah Sendiri
% <-- KTIDB : Karya Tulis Ilmiah dengan Bimbingan Tutor/Dosen
%=============================================================%

% --- Nama & Nomor Induk ---
\newcommand{\namaMahasiswa}{Yoeru Sandaru}
\newcommand{\nimMahasiswa}{081298765432}

% --- Program Studi ---
\newcommand{\prodi}{Sains Data}
\newcommand{\prodiENG}{Data Science} % <-- ENG

% --- Fakultas ---
\newcommand{\fakultas}{Sains dan Teknologi}
\newcommand{\fakultasENG}{Science and Technology} % <-- ENG

% --- Daerah & Negara ---
\newcommand{\daerahMahasiswa}{Medan}
\newcommand{\daerahKampus}{Medan}
\newcommand{\negaraMahasiswa}{Indonesia}

% --- Perguruan Tinggi ---
\newcommand{\kampus}{Universitas Terbuka}
\newcommand{\kampusENG}{Universitas Terbuka} % <-- ENG



%===============================%
%     IDENTITAS TUTOR/DOSEN     %
%===============================%=============================%
% Keterangan Fungsi Variabel
% [Kosong]  : Semua
% <-- ENG   : Dokumen Berbahasa Inggris
% <-- LJM   : Lembar Jawaban
% <-- BT    : Buku Tugas
% <-- KTI   : Karya Tulis Ilmiah Sendiri
% <-- KTIDB : Karya Tulis Ilmiah dengan Bimbingan Tutor/Dosen
%=============================================================%

% --- Nama & Nomor Induk ---
\newcommand{\namaPengampu}{Revo Wibowo, S.Pd., M.Pd., M.Sc.}
\newcommand{\namaPengampuPolos}{Revo Wibowo}